\documentclass{article}

% packages
\usepackage{amsmath, amsthm, thmtools, amsfonts, amssymb, luacode, catchfile, tikzducks, hyperref, ifthen}
\ifcsname c@kobocompile\endcsname
	\usepackage[a5paper, total={1072pt, 1448pt}, margin=10pt, includeheadfoot]{geometry} % set page margins
\else
	\usepackage[a4paper, margin=50pt, includeheadfoot]{geometry}
\fi
\usepackage[shortlabels]{enumitem}
\usepackage[skip=3pt, indent=0pt]{parskip}

% language
\usepackage[bidi=basic, layout=tabular, provide=*]{babel}
\ifcsname c@english\endcsname
	\babelprovide[main, import]{english}
\else
	\babelprovide[main, import]{hebrew}
	\babelprovide{rl}
\fi
%\babelfont{rm}{Libertinus Serif}
\babelfont{rm}[Renderer=Harfbuzz]{Libertinus Serif}
\babelfont{sf}{Libertinus Sans}
\babelfont{tt}{Libertinus Mono}

% style
\AddToHook{cmd/section/before}{\clearpage}	% Add line break before section
\linespread{1.3}
\setcounter{secnumdepth}{0}		% Remove default number tags from sections, this won't do well with theorems
\AtBeginDocument{\setlength{\belowdisplayskip}{3pt}}
\AtBeginDocument{\setlength{\abovedisplayskip}{3pt}}
\graphicspath{ {../images/} }

% operators
\DeclareMathOperator\cis{cis}
\DeclareMathOperator\Sp{Sp}
\DeclareMathOperator\tr{tr}
\DeclareMathOperator\im{Im}
\DeclareMathOperator\re{Re}
\DeclareMathOperator\diag{diag}
\DeclareMathOperator*\lowlim{\underline{lim}}
\DeclareMathOperator*\uplim{\overline{lim}}
\DeclareMathOperator\rng{rng}
\DeclareMathOperator\Sym{Sym}
\DeclareMathOperator\Arg{Arg}
\DeclareMathOperator\Log{Log}
\DeclareMathOperator\dom{dom}
\DeclareMathOperator\supp{Supp}
\DeclareMathOperator\var{Var}
\DeclareMathOperator\cov{Cov}

% commands
%\renewcommand\qedsymbol{\textbf{מש''ל}}
%\renewcommand\qedsymbol{\fbox{\emoji{lizard}}}
\newcommand{\Aa}[0]{\mathcal{A}}
\newcommand{\Bb}[0]{\mathcal{B}}
\newcommand{\CC}[0]{\mathbb{C}}
\newcommand{\Cc}[0]{\mathcal{C}}
\newcommand{\EE}[0]{\mathbb{E}}
\newcommand{\FF}[0]{\mathbb{F}}
\newcommand{\Ff}[0]{\mathcal{F}}
\newcommand{\Ii}[0]{\mathcal{I}}
\newcommand{\Gg}[0]{\mathcal{G}}
\newcommand{\Ll}[0]{\mathcal{L}}
\newcommand{\Mm}[0]{\mathcal{M}}
\newcommand{\NN}[0]{\mathbb{N}}
\newcommand{\Nn}[0]{\mathcal{N}}
\newcommand{\PP}[0]{\mathbb{P}}
\newcommand{\Pp}[0]{\mathcal{P}}
\newcommand{\QQ}[0]{\mathbb{Q}}
\newcommand{\RR}[0]{\mathbb{R}}
\newcommand{\Rr}[0]{\mathcal{R}}
\newcommand{\Ss}[0]{\mathcal{S}}
\newcommand{\TT}[0]{\mathbb{T}}
\newcommand{\Uu}[0]{\mathcal{U}}
\newcommand{\Vv}[0]{\mathcal{V}}
\newcommand{\Ww}[0]{\mathcal{W}}
\newcommand{\ZZ}[0]{\mathbb{Z}}
\newcommand{\acts}[0]{\circlearrowright}
\newcommand{\explain}[2] {
	\begin{flalign*}
		 && \text{#2} && \text{#1}
	\end{flalign*}
}
\newcommand{\maketitleprint}[0]{ \begin{center}
	%\begin{tikzpicture}[scale=3]
	%	\duck[graduate=gray!20!black, tassel=red!70!black]
	%\end{tikzpicture}	
	\includegraphics[width=6cm]{cover}
\end{center}
}

% theorem commands
\newtheoremstyle{c_remark}
	{}	% Space above
	{}	% Space below
	{}% Body font
	{}	% Indent amount
	{\bfseries}	% Theorem head font
	{}	% Punctuation after theorem head
	{.5em}	% Space after theorem head
	{\thmname{#1}\thmnumber{ #2}\thmnote{ \normalfont{\text{(#3)}}}}	% head content
\newtheoremstyle{c_definition}
	{3pt}	% Space above
	{3pt}	% Space below
	{}% Body font
	{}	% Indent amount
	{\bfseries}	% Theorem head font
	{}	% Punctuation after theorem head
	{.5em}	% Space after theorem head
	{\thmname{#1}\thmnumber{ #2}\thmnote{ \normalfont{\text{(#3)}}}}	% head content
\newtheoremstyle{c_plain}
	{3pt}	% Space above
	{3pt}	% Space below
	{\itshape}% Body font
	{}	% Indent amount
	{\bfseries}	% Theorem head font
	{}	% Punctuation after theorem head
	{.5em}	% Space after theorem head
	{\thmname{#1}\thmnumber{ #2}\thmnote{ \text{(#3)}}}	% head content

\ifcsname c@english\endcsname
	\theoremstyle{plain}
	\newtheorem{theorem}{Theorem}[section]
	\newtheorem{lemma}[theorem]{Lemma}
	\newtheorem{proposition}[theorem]{Proposition}
	\newtheorem*{proposition*}{Proposition}
	%\newtheorem{corollary}[theorem]{אין חלופה עברית}

	\theoremstyle{definition}
	\newtheorem{definition}[theorem]{Definition}
	\newtheorem*{definition*}{Definition}
	\newtheorem{example}{Example}[section]
	\newtheorem{exercise}{Exercise}[section]

	\theoremstyle{remark}
	\newtheorem*{remark}{Remark}
	\newtheorem*{solution}{Solution}
	\newtheorem{conclusion}[theorem]{Conclusion}
	\newtheorem{notation}[theorem]{Notation}
\else
	\theoremstyle{c_plain}
	\newtheorem{theorem}{משפט}[section]
	\newtheorem{lemma}[theorem]{למה}
	\newtheorem{proposition}[theorem]{טענה}
	\newtheorem*{proposition*}{טענה}
	%\newtheorem{corollary}[theorem]{אין חלופה עברית}

	\theoremstyle{c_definition}
	\newtheorem{definition}[theorem]{הגדרה}
	\newtheorem*{definition*}{הגדרה}
	\newtheorem{example}{דוגמה}[section]
	\newtheorem{exercise}{תרגיל}[section]

	\theoremstyle{c_remark}
	\newtheorem*{remark}{הערה}
	\newtheorem*{solution}{פתרון}
	\newtheorem{conclusion}[theorem]{מסקנה}
	\newtheorem{notation}[theorem]{סימון}
\fi

% Questions related commands
\newcounter{question}
\setcounter{question}{1}
\newcounter{sub_question}
\setcounter{sub_question}{1}

\ifcsname c@english\endcsname
	\newcommand{\question}[1][0]{
		\ifthenelse{#1 = 0}{}{\setcounter{question}{#1}}
		\section{Question \arabic{question}}
		\addtocounter{question}{1}
		\setcounter{sub_question}{1}
	}

	\newcommand{\subquestion}[1][0]{
		\ifthenelse{#1 = 0}{}{\setcounter{sub_question}{#1}}
		\subsection{Part \alph{sub_question}}
		\addtocounter{sub_question}{1}
	}
\else
	\newcommand{\question}[1][0]{
		\ifthenelse{#1 = 0}{}{\setcounter{question}{#1}}
		\section{שאלה \arabic{question}}
		\addtocounter{question}{1}
		\setcounter{sub_question}{1}
	}

	\newcommand{\subquestion}[1][0]{
		\ifthenelse{#1 = 0}{}{\setcounter{sub_question}{#1}}
		\subsection{סעיף \localecounter{letters.gershayim}{sub_question}}
		\addtocounter{sub_question}{1}
	}
\fi

% import lua and start of document
\directlua{common = require ('../common')}

\GetEnv{AUTHOR}

% headers
\author{\AUTHOR}
\date\today


\title{פתרון ממ''ן 11 – אלגברה לינארית 2 (20229)}

\begin{document}
\maketitle

\section{שאלה 1}
\subsection{סעיף א'}
יהיה $V = M_{n \times n}^\CC$ עם המכפלה הפנימית הסטנדרטית, תהי $P \in V$ מטריצה הפיכה,
ותהי העתקה לינארית $T_P: V \to V$ המוגדרת
\[
	T_P X = P^{-1} X P
\]
נוכיח שמתקיים $T_P^* = T_{P^*}$.
\begin{proof}
	נגדיר $A, B \in V$. נשים לב כי לפי משפט 2.1.4(ו') מתקיים
	\[
		{(P^{-1})}^* = {(P^*)}^{-1} \tag{1}
	\]
	נראה כי
	\begin{align*}
		(A, T_{P^*} B)
		& = (A, {(P^*)}^{-1} B P^*) \\
		& = \tr( {({(P^*)}^{-1} B P^*)}^* A ) \\
		& = \tr( {(P^*)}^* {({(P^*)}^{-1} B)}^* A ) && \text{2.1.4 ו'} \\
		& = \tr( {(P^*)}^* B^* {({(P^*)}^{-1})}^* A ) \\
		& = \tr( {(P^*)}^* B^* {({(P^*)}^{-1})}^* A ) \\
		& = \tr( P B^* P^{-1} A ) && (1) \\
		& = \tr( B^* P^{-1} A P ) && \text{מכפלה ב־tr} \\
		& = (P^{-1} A P, B ) \\
		& = (T_P A, B ) \\
	\end{align*}
	על־פי הגדרת העתקה צמודה מתקיים
	\[
		{(T_P)}^* = (T_{P^*})
	\]
\end{proof}

\subsection{סעיף ב'}
נגדיר $V = M^\CC_{2 \times 2}$.
תהי $T_P : V \to V$ המוגדרת על־ידי $T_P X = P^{-1} X P$ כאשר
\[
	P = \begin{pmatrix}
		i & 1 \\
		-1 & -i
	\end{pmatrix}
\]
נמצא את המטריצה המייצגת את ${(T_P)}^*$ בבסיס הסטנדרטי של $V$. \\*
על־פי הסעיף הקודם מתקיים ${(T_P)}^* = T_{P^*}$, ולפי חישוב
\[
	P^* = \begin{pmatrix}
		-i & -1 \\
		1 & i
	\end{pmatrix},
	{(P^*)}^{-1} = \frac{1}{2}\begin{pmatrix}
		i & 1 \\
		-1 & -i
	\end{pmatrix}
\]
נגדיר $B = (E_1, E_2, E_3, E_4)$ הבסיס הסטנדרטי של $V$ ונחשב:
\[
	T_{P^*} E_1 = {(P^*)}^{-1} E_1 P^*
	= \frac{1}{2}
	\begin{pmatrix} -i & -1 \\ 1 & i \end{pmatrix}
	\begin{pmatrix} 1 & 0 \\ 0 & 0 \end{pmatrix}
	\begin{pmatrix} i & 1 \\ -1 & -i \end{pmatrix}
	= \frac{1}{2}
	\begin{pmatrix} -i & 0 \\ 1 & 0 \end{pmatrix}
	\begin{pmatrix} i & 1 \\ -1 & -i \end{pmatrix}
	= \frac{1}{2} \begin{pmatrix} 1 & -i \\ i & 1 \end{pmatrix}
\]
\[
	T_{P^*} E_2 = \frac{1}{2}
	\begin{pmatrix} -i & -1 \\ 1 & i \end{pmatrix}
	\begin{pmatrix} 0 & 1 \\ 0 & 0 \end{pmatrix}
	\begin{pmatrix} i & 1 \\ -1 & -i \end{pmatrix}
	= \frac{1}{2}
	\begin{pmatrix} 0 & -i \\ 0 & 1 \end{pmatrix}
	\begin{pmatrix} i & 1 \\ -1 & -i \end{pmatrix}
	= \frac{1}{2} \begin{pmatrix} i & -1 \\ -1 & -i \end{pmatrix}
\]
\[
	T_{P^*} E_3 = \frac{1}{2}
	\begin{pmatrix} -i & -1 \\ 1 & i \end{pmatrix}
	\begin{pmatrix} 0 & 0 \\ 1 & 0 \end{pmatrix}
	\begin{pmatrix} i & 1 \\ -1 & -i \end{pmatrix}
	= \frac{1}{2}
	\begin{pmatrix} -1 & 0 \\ i & 0 \end{pmatrix}
	\begin{pmatrix} i & 1 \\ -1 & -i \end{pmatrix}
	= \frac{1}{2} \begin{pmatrix} -i & -1 \\ -1 & i \end{pmatrix}
\]
\[
	T_{P^*} E_4 = \frac{1}{2}
	\begin{pmatrix} -i & -1 \\ 1 & i \end{pmatrix}
	\begin{pmatrix} 0 & 0 \\ 0 & 1 \end{pmatrix}
	\begin{pmatrix} i & 1 \\ -1 & -i \end{pmatrix}
	= \frac{1}{2}
	\begin{pmatrix} 0 & -1 \\ 0 & i \end{pmatrix}
	\begin{pmatrix} i & 1 \\ -1 & -i \end{pmatrix}
	= \frac{1}{2} \begin{pmatrix} 1 & i \\ -i & 1 \end{pmatrix}
\]
לכן 
\[
	{[T_{P^*}]}_B = \begin{pmatrix}
		1 & -i & i & 1 \\
		i & -1 & -1 & -i \\
		-i & -1 & -1 & i \\
		1 & i & -i & 1
	\end{pmatrix}
\]

\section{שאלה 2}
\subsection{סעיף א'}
יהיו $P, Q \in M^\RR_{n \times n}$ ותהי $U = P + iQ$. נסמן
\[
	D = \begin{pmatrix}
		P & -Q \\
		Q & P
	\end{pmatrix}
\]
נוכיח שאם $U$ מטריצה הרמיטית אז $D$ מטריצה סימטרית.
\begin{proof}
	נניח כי $U$ הרמיטית, לכן מהגדרה 1.2.4 וממשפט 2.1.4 נובע כי:
	\[
		U^* = {(P + iQ)}^* = P^* - i Q^* = P + i Q
	\]
	בשל היות $P, Q$ ממשיות מתקיים:
	\[
		P^* = P^t = P, -Q^* = -Q^t = Q
	\]
	לכן $P$ מטריצה סימטרית ו־$Q$ מטריצה אינטי־סימטרית, ובהתאם:
	\[
		D^t = \begin{pmatrix}
			P^t & Q^t \\
			{(-Q)}^t & P^t
		\end{pmatrix}
		= \begin{pmatrix}
			P & -Q \\
			Q & P
		\end{pmatrix}
		= D
	\]
	המטריצה $D$ היא סימטרית.
\end{proof}

\subsection{סעיף ב'}
נוכיח כי אם $U$ מטריצה אוניטרית אז $D$ מטריצה אורתוגונלית.
\begin{proof}
	נניח כי $U$ אוניטרית, לכן על־פי הגדרה 2.3.1
	\[
		U U^* = I \rightarrow (P + i Q)(P^t - i Q^t) = P P^t + i Q P^t - i P Q^t + Q Q^t = I \rightarrow PP^t + QQ^t = I, QP^t = PQ^t
	\]
	נחשב
	\[
		D D^t
		= \begin{pmatrix}
			P & -Q \\
			Q & P
		\end{pmatrix} \begin{pmatrix}
			P^t & Q^t \\
			-Q^t & P^t
		\end{pmatrix}
		= \begin{pmatrix}
			PP^t + QQ^t & PQ^t - QP^t \\
			QP^t - PQ^t & QQ^t + PP^t
		\end{pmatrix}
		= \begin{pmatrix}
			I_n & 0 \\
			0 & I_n
		\end{pmatrix}
		= I_{2n}
	\]
	לכן $D$ מטריצה אורתוגונלית.
\end{proof}

\section{שאלה 3}
תהינה $A, B$ מטריצות חיוביות לחלוטין ו־$Q$ מטריצה אוניטרית.
\begin{proof}
	נוכיח כי אם $A = BQ$ אז $A = B$. \\*
	יהי $\lambda$ ערך עצמי של $Q$ ו־$u$ וקטור עצמי עבור $\lambda$, לכן $A u = B Q u = \lambda B u$. \\*
	נשים לב כי
	\[
		(A u, u) = (B Q u, u) = \lambda (B u, u)
	\]
	ידוע כי $A, B$ חיוביות לחלוטין, לכן $(Au, u), (Bu, u) > 0$, ובהתאם גם
	\[
		\lambda = \frac{(A u, u)}{(B u, u)} > 0
	\]
	לכן בכלל גם $\lambda \in \RR$, אז לפי טענה 2.4.3 ואי־השוויון מתקיים $\lambda = 1$ וזהו הערך העצמי היחיד של המטריצה. \\*
	$Q$ מטריצה אוניטרית ולכן בכלל גם נורמלית כמסקנה מהגדרה 3.1.4, לכן ממשפט 3.1.4 נובע כי $Q$ דומה למטריצה אלכסונית.
	עוד ידוע כי ל־$Q$ ערך עצמי יחיד $1$ ולכן $Q$ דומה למטריצה אלכסונית שאלכסונה הוא $1$, דהינו למטריצת היחידה.
	לכן קיימת מטריצה הפיכה $P$ כך שמתקיים $Q = P^{-1} I P = P^{-1} P = I$ וכן $Q = I$. \\*
	אז גם $A = B Q = B I = B$.
\end{proof}

\section{שאלה 4}
יהי $0 \ne w \in \CC^n$ וקטור עמודה. נמצא תנאי הכרחי ומספיק עבור $w$ כדי שהמטריצה $H = I - 2 w w^*$ תהיה אוניטרית. \\*
נשים לב כי $H^* = {(I - 2 ww^*)}^* = I^* - {(2ww^*)}^* = I - 2{(ww^*)}^*$. \\*
ידוע גם כי $w^* = \overline{w}^t$, וגם ${\overline{(w w^*)}}^t = {(\overline{w} w^t)}^t = w w^*$ לפי חוקי שחלוף מהקורס הקודם. \\*
על־כן מתקיים $H^* = H$. נחשב:
\begin{align*}
	H H^* & = H^* H = H^2 = (I - 2w w^*)(I - 2w w^*) \\
		  & = I^2 - 2 \cdot I \cdot 2w w^* + 4{(ww^*)}^2 = I - 4ww^* + 4w \lVert w \rVert w^* \\
		  & = I - 4 ww^* (1 - \lVert w \rVert)
\end{align*}
אנו יודעים כי $w \ne 0$ ולכן גם $ww^* \ne 0$, בהתאם מתקיים $H H^* = I$ כאשר $\lVert w \rVert = 1$ ולכן זהו תנאי הכרחי ומספיק לאוניטריות $H$. \\*
נשים לב כי במקרה זה $H$ היא מטריצת שיקוף ביחס ל־${\{ w \}}^\perp$, שכן על־פי חישוב ישיר מתקיים:
\[
	H w = I w - 2w w^* w = w - 2 \lVert w \rVert w = w - 2w = -w
\]
\[
	\forall v \in {\{ w \}}^\perp: H v = I v - 2 w w^* v = v - 2w 0 = v
\]

\section{שאלה 5}
יהי $V$ מרחב מכפלה פנימית, $w_1, w_2 \in V$ וקטורים המקיימים
\[
	(w_1, w_2) = 0, \lVert w_1 \rVert = \lVert w_2 \rVert = 1
\]
נגדיר העתקה לינארית $T: V \to V$ כך שמתקיים
\[
	Tv = v - 2(v, w_1) w_1 - 2(v, w_2) w_2
\]

\subsection{סעיף א'}
נוכיח כי $T$ צמודה לעצמה ואוניטרית.
\begin{proof}
	על־פי למה 1.2.3:
	\begin{align*}
		(T v, u)
		& = (v - 2(v, w_1) w_1 - 2(2, w_2) w_2, u) \\
		& = (v, u) - 2 (v, w_1) (w_1, u) - 2(v, w_2) (w_2, u) \\
		& = (v, u) - (v, 2 (u, w_1) w_1) - (v, (u, w_2) 2w_2) \\
		& = (v, u) - (v, 2 (u, w_1) w_1) - (v, (u, w_2) 2w_2) \\
		& = (v, u - 2 (u, w_1) w_1 - (u, w_2) 2w_2) \\
		(T v, u) & = (v, T u)
	\end{align*}
	לכן $T$ צמודה לעצמה. \*
	נחשב
	\begin{align*}
		(T v, T v)
		= & (v - 2(v, w_1) w_1 - 2(2, w_2) w_2, v - 2(v, w_1) w_1 - 2(2, w_2) w_2) \\
		= & (v, v - 2(v, w_1) w_1 - 2(2, w_2) w_2) \\
		& - ((v, w_1) w_1, v - 2(v, w_1) w_1 - 2(2, w_2) w_2) \\
		& - 2((v, w_2) w_2, v - 2(v, w_1) w_1 - 2(2, w_2) w_2) \\
		= & (v, v) - 2 {(v, w_1)}^2 - 2{(v, w_2)}^2 \\
		& - 2{(v, w_1)}^2 + 4{(v, w_1)}^2  + 4{(v, w_2)}^2 \\
		& - 2{(v, w_2)}^2 + 4{(v, w_1)}^2 + 4{(2, w_2)}^2 \\
		= & \lVert v\rVert^2
	\end{align*}
	ובהתאם להגדרת הנורמה $\lVert T v \rVert = \lVert v \rVert$ ולכן לפי משפט 2.3.2 $T$ אוניטרית.
\end{proof}

\subsection{סעיף ב'}
נבדוק אם $T$ אי־שלילית. \\*
נחשב את $(T w_1, w_1)$:
\[
	(T w_1, w_1)
	= (w_1, w_1) - 2 {(w_1, w_1)}^2 - 2{(w_1, w_2)}^2
	= \lVert w_1 \rVert - 2\lVert w_1 \rVert^2 - 0
	= -1
\]
לכן מהגדרה 2.2.7 נובע ישירות כי $T$ איננה אי־שלילית.

\end{document}
